\documentclass{article}


% if you need to pass options to natbib, use, e.g.:
%     \PassOptionsToPackage{numbers, compress}{natbib}
% before loading neurips_2025


% ready for submission
% \usepackage{neurips_2025}


% to compile a preprint version, e.g., for submission to arXiv, add add the
% [preprint] option:
\usepackage[preprint]{neurips_2025}


% to compile a camera-ready version, add the [final] option, e.g.:
%     \usepackage[final]{neurips_2025}


% to avoid loading the natbib package, add option nonatbib:
%    \usepackage[nonatbib]{neurips_2025}


\usepackage[utf8]{inputenc} % allow utf-8 input
\usepackage[T1]{fontenc}    % use 8-bit T1 fonts
\usepackage{hyperref}       % hyperlinks
\usepackage{url}            % simple URL typesetting
\usepackage{booktabs}       % professional-quality tables
\usepackage{tabularx}   % For tables with adjustable-width columns
\usepackage{siunitx} 
\usepackage{amsfonts}       % blackboard math symbols
\usepackage{amsmath}       % blackboard math symbols
\usepackage{nicefrac}       % compact symbols for 1/2, etc.
\usepackage{microtype}      % microtypography
\usepackage{xcolor}         % colors
\usepackage{graphicx} 
\usepackage{subcaption}
\usepackage{epsfig}
\usepackage{array}     
\usepackage{multirow}
\usepackage{makecell} 
\usepackage{hhline}
\usepackage{enumitem}
\usepackage[dvipsnames]{xcolor}


\newcommand{\LHC}[1]{{\color{blue}#1}}

% \title{UniVerse-1: A Unified Audio-Video Synthesis Model via Model-Stitching}
\title{UniVerse-1: \underline{Uni}fied Audio-\underline{V}ideo Gen\underline{er}ation via \underline{S}titching of \underline{E}xperts}
% From Snapshot to Spotlight: Multi-Modal Driven Avatar Generation from a Human Image


% The \author macro works with any number of authors. There are two commands
% used to separate the names and addresses of multiple authors: \And and \AND.
%
% Using \And between authors leaves it to LaTeX to determine where to break the
% lines. Using \AND forces a line break at that point. So, if LaTeX puts 3 of 4
% authors names on the first line, and the last on the second line, try using
% \AND instead of \And before the third author name.


% \author{%
%   David S.~Hippocampus\thanks{Use footnote for providing further information
%     about author (webpage, alternative address)---\emph{not} for acknowledging
%     funding agencies.} \\
%   Department of Computer Science\\
%   Cranberry-Lemon University\\
%   Pittsburgh, PA 15213 \\
%   \texttt{hippo@cs.cranberry-lemon.edu} \\
% }
\author{
    % \textbf{Duomin Wang}\thanks{Equal contribution}$^{*1}$ \quad 
    \textbf{Duomin Wang}$^{1}$ \quad 
    \textbf{Wei Zuo}$^{1}$ \quad 
    \textbf{Aojie Li}$^{1}$ \quad 
    \textbf{Ling-Hao Chen}$^{1,4}$ \quad
    \textbf{Xinyao Liao}$^{1}$ \quad  \\
    \textbf{Deyu Zhou}$^{1,2}$ \quad 
    \textbf{Zixin Yin}$^{1,3}$ \quad 
    \textbf{Xili Dai}$^{2}$ \quad 
    \textbf{Daxin Jiang}$^{1}$ \quad 
    \textbf{Gang Yu}$^{1}$ \quad 
    \\
    StepFun$^{1}$ \quad \\
    The Hong Kong University of Science and Technology(GuangZhou)$^{2}$ \quad \\
    The Hong Kong University of Science and Technology$^{3}$ Tsinghua University$^{4}$\quad \\
}

% If you wish to post a preprint of your work online, e.g., on arXiv, using the NeurIPS style, please use the \verb+preprint+ option. This will create a nonanonymized version of your work with the text ``Preprint. Work in progress.'' in the footer. This version may be distributed as you see fit, as long as you do not say which conference it was submitted to. Please \textbf{do not} use the \verb+final+ option, which should \textbf{only} be used for papers accepted to NeurIPS.
\begin{document}


\maketitle


\begin{abstract}
We introduce \textbf{UniVerse-1}, a unified, Veo-3-like model capable of simultaneously generating coordinated audio and video. To enhance training efficiency, we bypass training from scratch and instead employ a \textbf{stitching of experts (SoE)} technique. This approach deeply fuses the corresponding blocks of pre-trained video and music generation experts models, thereby fully leveraging their foundational capabilities. To ensure accurate annotations and temporal alignment for both ambient sounds and speech with video content, we developed an \textbf{online annotation pipeline} that processes the required training data and generates labels during training process. This strategy circumvents the performance degradation often caused by misalignment text-based annotations. Through the synergy of these techniques, our model, after being finetuned on approximately $7,600$ hours of audio-video data, produces results with well-coordinated audio-visuals for ambient sounds generation and strong alignment for speech generation. To systematically evaluate our proposed method, we introduce \textbf{Verse-Bench}, a new benchmark dataset. In an effort to advance research in audio-video generation and to close the performance gap with state-of-the-art models such as Veo3, we make our model and code publicly available. We hope this contribution will benefit the broader research community. Project page: \url{https://dorniwang.github.io/UniVerse-1/}.
\end{abstract}

% \section{Introduction}

% The era of diffusion models~\cite{song2020score, song2020improved, lipman2022flow} has heralded a period of monumental development in video generation. The trajectory began with foundational UNet-based diffusion models~\cite{guo2023animatediff,singer2022make,ho2022video,xing2024dynamicrafter} and has recently culminated in DiT-based architectures~\cite{kong2024hunyuanvideo,yang2024cogvideox,wan2025wan,polyak2024movie,ma2025step}, pioneered by landmark models like Sora~\cite{openaisora2024}. This new generation of video models has achieved unprecedented progress in both prompt alignment and generation quality. These exhilarating advancements are catalyzing a profound transformation in the field of video creation, unlocking limitless possibilities.
% Commercial applications, such as film production, advertising, and creative short-form video generation, stand to benefit immensely. By leveraging video generation technology, professionals can significantly enhance their workflow efficiency and reduce production costs. Concurrently, the evolution of these powerful video models has spurred a wave of research into downstream applications built upon pre-trained backbones. Fields such as talking head synthesis~\cite{wang2023progressive,yu2023talking,wang2025fantasytalking,luo2025dreamactor}, video-based virtual try-on~\cite{ning2024picture,wang2025mv,kim2024stableviton}, human animation~\cite{tan2024animate,chen2025hunyuanvideo,lin2025omnihuman}, and video editing~\cite{xue2025stand,jiang2025vace} have all seen remarkable breakthroughs, empowered by the enhanced capabilities of modern video models. These applications are now providing viable solutions for a multitude of real-world scenarios.

% The current generation of video synthesis, driven by diffusion foundation models, has evolved along several key axes. A primary development has been in model architecture, marked by a transition from early UNet-based frameworks~\cite{rombach2022high} to the now-prevalent Diffusion Transformer (DiT) architecture~\cite{peebles2023scalable}. Initial models, such as AnimateDiff~\cite{guo2023animatediff} and Make-A-Video~\cite{singer2022make}, extended pre-trained UNet image generators by incorporating temporal layers to enforce inter-frame consistency. A paradigm shift occurred with the advent of Sora~\cite{openaisora2024}, followed by open-source DiT models like CogVideox~\cite{yang2024cogvideox}, HunyuanVideo~\cite{kong2024hunyuanvideo}, and WAN2.1~\cite{wan2025wan}. These models, trained from scratch with 3D full-attention mechanisms, have demonstrated significantly more promising results.

% Furthermore, the success of Sora has underscored the importance of scaling laws for both data and model size. A key insight for the research community is that systematically scaling up a DiT-based architecture leads to substantial improvements in video generation performance. This principle has been validated not only by open-source efforts but also by a series of high-performing, closed-source models such as SeeDance1.0~\cite{gao2025seedance}, Kling~\cite{kuaishou2024kling}, Veo2~\cite{veo22025}, etc., all of which leverage the DiT architecture at a massive scale.

% Architecturally, these models share a common blueprint: a 3D Variational Autoencoder (VAE) first compresses the video into a latent representation. The DiT then operates on this latent space, executing a conditional denoising process on noisy latents to reverse the diffusion and generate the final video. The development efforts for these models have predominantly focused on two objectives: enhancing semantic alignment to text prompts and improving the overall visual quality of the output. To this end, researchers have employed a suite of techniques, including curating high-quality datasets, optimizing video captioning models, advanced prompt engineering, integrating super-resolution modules, and implementing post-training alignment strategies such as Supervised Fine-Tuning (SFT)~\cite{instag} and Reinforcement Learning from Human Feedback (RLHF)~\cite{rlhf}. These validated methods have continuously advanced model capabilities and, in turn, have empowered content creators by streamlining the initial video production workflow, thereby reducing time and cost while boosting creative efficiency.

% Despite these advancements, a significant limitation persists: a near-exclusive focus on the visual modality. Video is an inherently multimodal medium, combining both visual and audio streams. While post-hoc video-to-audio models can generate sound for synthesized videos, this decoupled approach is fundamentally flawed. It can adapt the audio to the video content, but it cannot enforce temporal alignment in the reverse direction. For example, it is incapable of synchronizing the lip movements of a speaking character with the corresponding speech. Google's Veo3~\cite{veo32025} has demonstrated the ability to generate audio and video synchronously, but it remains a closed-source model with no publicly available technical details.

% To address this critical gap between closed-source commercial systems and open research, we introduce \textbf{UniVerse-1}: a unified, fully open-source, Veo-3-like model capable of simultaneously generating coordinated audio and video. 

% To enhance training efficiency and capitalize on the robust capabilities of existing pre-trained models, our approach circumvents training a new model from scratch. Instead, we propose a \textbf{stitching of experts (SoE)} methodology that integrates a state-of-the-art video generation model, WAN2.1~\cite{wan2025wan}, with a music generation model, Ace-step~\cite{gong2025ace}. The integration of two experts is realized by introducing cross-modal MLP connectors within each corresponding block of the two models, enabling bidirectional interaction between the video and audio modalities. This strategy was found to significantly accelerate training convergence.
% To preserve the critical temporal and semantic alignment between audio-video data and their textual descriptions, we developed an \textbf{online annotation pipeline}. This pipeline generates labels dynamically during the training process, thereby mitigating the performance degradation that can arise from static, misaligned annotations.
% Furthermore, our investigation revealed a crucial insight for bimodal diffusion modeling: the necessity of independent noise sampling for each modality. We identified that without such isolation, the pseudo-random number generation process~\cite{hamming1952mathematical} can introduce spurious correlations between the noise vectors for video and audio, which subsequently degrades audio quality during inference. To support this work, we curated a high-quality dataset comprising approximately $7,600$ hours of precisely aligned audio-video content. To systematically evaluate our proposed method, we introduce \textbf{Verse-Bench}, a new benchmark dataset. Verse-Bench comprises $600$ image-text prompt pairs designed to cover a diverse range of sound categories. Highlighting its versatility, the benchmark supports the evaluation of not only joint audio-video generation but also unidirectional tasks. It can be used to assess video-to-audio generation, and it includes the specialized Verse-Ted subset for evaluating audio-to-video tasks like talking head synthesis.

% In summary, our primary contributions are:
% \begin{itemize}[leftmargin=*]
%     \item An Open-source Audio-Video Foundation Model: We present a novel, open-source model for synchronous audio-video generation that produces highly coherent and well-aligned audio-visual content.
%     \item A Pre-trained Model Fusion Strategy: We introduce a stitching of experts technique that effectively fuses the capabilities of existing pre-trained models, significantly accelerating convergence while leveraging their powerful priors.
%     \item An Online Data Annotation Pipeline: We developed an online data annotation pipeline that ensures strict temporal and semantic alignment between training data and captions, a crucial factor for high-fidelity bimodal generation.
%     \item A new Verse-Bench is proposed to comprehensively evaluate joint audio-video generation models.
% \end{itemize}


\section{Introduction}

The era of diffusion models~\cite{song2020score, song2020improved, lipman2022flow} has culminated in the rise of Diffusion Transformer (DiT) architectures~\cite{peebles2023scalable}, exemplified by landmark models like Sora~\cite{openaisora2024} and its open-source counterparts~\cite{kong2024hunyuanvideo,yang2024cogvideox,wan2025wan}. These models, leveraging unprecedented scales of data and computation, have achieved remarkable quality and prompt alignment in video generation. This success is catalyzing a profound transformation across creative industries and has spurred a wave of research into downstream applications such as talking head synthesis~\cite{wang2023progressive,yu2023talking,wang2025fantasytalking,luo2025dreamactor} and human animation~\cite{tan2024animate,chen2025hunyuanvideo,lin2025omnihuman}, which now offer viable real-world solutions.

However, this rapid progress has been almost exclusively confined to the visual domain, treating video as a silent movie. This \textbf{unimodal focus} represents a fundamental bottleneck, as it ignores the inherently multimodal nature of video. Post-hoc video-to-audio models~\cite{shan2025hunyuanvideo} serve as a superficial fix, but they are inherently limited; while capable of adapting audio to existing visual content, they fail to enforce temporal alignment in the reverse direction. This makes critical tasks, such as synchronizing lip movements with speech, impossible. While closed-source systems like Google's Veo3~\cite{veo32025} have demonstrated synchronous audio-video generation, the lack of publicly available technical details leaves a critical gap in open research.

To bridge this gap between closed-source systems and open research, we introduce \textbf{UniVerse-1}: a unified, fully open-source, Veo-3-like model capable of simultaneously generating coordinated audio and video. Our work is underpinned by several key technical contributions designed to address the unique challenges of bimodal generation.

Instead of the costly process of training a new model from scratch, we propose a novel and efficient \textbf{stitching of experts (SoE)} paradigm. This methodology effectively fuses a state-of-the-art video generation model, WAN2.1~\cite{wan2025wan}, with a music generation model, Ace-step~\cite{gong2025ace}. The core of this fusion lies in lightweight, cross-modal MLP connectors introduced within corresponding blocks of each model. These connectors facilitate bidirectional interaction between modalities, and we found this strategy to significantly accelerate training convergence by leveraging the powerful priors of the pre-trained experts.

Furthermore, we tackle the critical challenge of data alignment in bimodal training. We argue that static, pre-processed annotations are a flawed paradigm for tasks requiring precise temporal consistency. To address this, we developed an \textbf{online annotation pipeline} that generates labels dynamically during training. This approach ensures strict temporal and semantic alignment between audio-video data and their textual descriptions, mitigating the performance degradation caused by static misalignment. During our investigation, we also uncovered a crucial, yet overlooked, factor in bimodal diffusion modeling: \textbf{cross-modal noise correlation}. We identified that the standard pseudo-random number generation process~\cite{hamming1952mathematical} can introduce spurious correlations between the noise vectors for video and audio, which subsequently degrades audio quality during inference. Our solution involves ensuring independent noise sampling for each modality.

To support this work, we curated a high-quality dataset comprising approximately $7,600$ hours of precisely aligned audio-video content. To systematically evaluate our method, we also propose \textbf{Verse-Bench}, a new benchmark featuring $600$ image-text prompt pairs covering a diverse range of sound categories. Highlighting its versatility, Verse-Bench supports not only joint audio-video generation but also unidirectional tasks, including a specialized \textbf{Verse-Ted} subset designed for evaluating audio-to-video synthesis.

In summary, our primary contributions are:
\begin{itemize}[leftmargin=*]
    \item \textbf{An Open-source Audio-Video Foundation Model:} We present UniVerse-1, a novel, open-source model capable of producing highly coherent and well-aligned synchronous audio-visual content, closing a critical gap in the open-source community.
    
    \item \textbf{A Novel Methodology for Joint Audio-Video Generation:} We propose a comprehensive methodology to enable efficient and high-quality joint audio-video synthesis. This is achieved through three key innovations: a \textbf{stitching of experts (SoE)} paradigm to accelerate convergence by fusing pre-trained models; an \textbf{online data annotation pipeline} to solve the critical static misalignment problem in training; and the identification and mitigation of the previously overlooked \textbf{cross-modal noise correlation} issue, a crucial factor for generation quality.
    
    \item \textbf{A Comprehensive Evaluation Benchmark:} We propose Verse-Bench, a new benchmark designed to comprehensively evaluate joint audio-video generation models across a diverse set of tasks.
\end{itemize}
\section{Related Works}
\paragraph{Video Diffusion Models} The field of video generation was revolutionized by the introduction of diffusion models, with pioneering works like AnimateDiff~\cite{guo2023animatediff} and Video Diffusion Models~\cite{ho2022video} marking the beginning of this new era. This initial wave of research was further advanced by models such as Stable Video Diffusion~\cite{blattmann2023stable}, which first demonstrated that curating large-scale, high-quality datasets is critical for enhancing model performance. A common characteristic of these early models was their reliance on UNet architectures. To mitigate the challenges posed by the limited availability and quality of video data compared to images, these models were typically fine-tuned from pre-trained UNet image foundations.
A significant paradigm shift occurred with the introduction of Sora~\cite{openaisora2024}, which heralded a new age defined by the Diffusion Transformer (DiT) architecture~\cite{peebles2023scalable} and training on massive, high-quality video corpora. This breakthrough spurred a proliferation of subsequent research. CogVideox~\cite{yang2024cogvideox} was the first to release an open-source DiT-based model, providing a significant catalyst for community-driven innovation. This was followed by other notable open-source models such as HunyuanVideo~\cite{kong2024hunyuanvideo}, WAN2.1~\cite{wan2025wan}, and Step-Video~\cite{ma2025step}, as well as high-performing closed-source systems including Kling~\cite{kuaishou2024kling}, SeeDance 1.0~\cite{gao2025seedance}, Movie Gen~\cite{polyak2024movie}, and Veo2~\cite{veo22025}.
Architecturally, these contemporary models converge on a common blueprint. They employ a 3D Variational Autoencoder (VAE) to achieve spatio-temporal compression of video into a latent space. The core generative process is then handled by a DiT, which learns to denoise these noisy latents. Across these state-of-the-art models, the quality and scale of the training data have been identified as paramount factors, making data curation and processing a central component of their development.

\paragraph{Audio Diffusion Models} The application of diffusion models to audio generation has followed a parallel trajectory to their video counterparts, fundamentally transforming the landscape of text-to-audio and text-to-music synthesis. Early explorations demonstrated the potential of diffusion for generating high-fidelity audio, but a pivotal advancement was the adoption of latent diffusion architectures~\cite{rombach2022high}, which significantly improved both efficiency and quality. A common technical pipeline for these approaches involves first transforming the raw audio waveform into a mel spectrogram. A Variational Autoencoder (VAE) is then trained on this spectrogram representation to learn a compressed latent space, within which the core diffusion process operates. Within this framework, models such as Stable Audio Open~\cite{evans2025stable} and Riffusion~\cite{riffusion} excel at generating high-fidelity, long-form audio. Furthermore, models like the AudioLDM series~\cite{audioldm2-2024taslp,liu2023audioldm} and DiffRhythm~\cite{ning2025diffrhythm} advance these capabilities to vocal music synthesis, offering fine-grained control over rhythm and other expressive attributes.

\paragraph{Joint Audio and Video Generation} The exploration of joint audio-video generation within diffusion frameworks began with pioneering efforts like MM-Diffusion~\cite{ruan2023mm}. This model was the first to tackle this bimodal task, employing a UNet architecture with two distinct subnetworks, each dedicated to processing the audio and video modalities, respectively. Following this initial work, a series of subsequent models emerged~\cite{ishii2024simple,hayakawa2024mmdisco,ergasti2025r}. However, these early approaches were typically constrained by small-scale training datasets~\cite{lee2022sound,li2021ai}, often less than $10$ hours in size, which inherently limited their diversity and generalization capabilities. A notable step forward was made by models such as Syncflow~\cite{liu2024syncflow} and Uniform~\cite{zhao2025uniform}, which scaled up the training data to approximately $500$ hours by leveraging the VGGSound~\cite{chen2020vggsound} and AudioSet~\cite{gemmeke2017audio} dataset, thereby enhancing their generalization. Despite this progress, persistent challenges remained, including suboptimal video quality and a low degree of disentanglement between the audio and visual streams. The advent of Google's Veo3~\cite{veo32025} marked a significant milestone, representing the first large-scale initiative in synchronous audio-video generation. Veo3 demonstrated the capacity for generating high-fidelity audio and video that is not only diverse but also semantically and temporally coordinated, strictly adhering to user-provided text prompts.
\section{UniVerse-1}

\subsection{Preliminary}
\label{sec:preliminary}

% Our model is constructed upon the foundations of the Wan2.1 (1.3B parameters)~\cite{wan2025wan} text-to-video model and the Ace-step (3.5B parameters)~\cite{gong2025ace} music generation model. We will briefly introduce their respective architectures.

% \LHC{
Our model is constructed upon the foundations of the Wan2.1 (1.3B parameters)~\cite{wan2025wan} text-to-video model and the Ace-step (3.5B parameters)~\cite{gong2025ace} music generation model. Before delving into technical details, we will briefly introduce their respective architectures.
% }


% Wan2.1 is composed of three primary components: a 3D Variational Autoencoder (VAE), an umT5~\cite{chung2023unimax} text encoder, and a Diffusion Transformer (DiT). The 3D VAE compresses an input video of shape $(3, T, H, W)$ into a latent representation of shape $(16, T/t, H/h, W/w)$, where the temporal and spatial downsampling factors are $t=4$ and $h=w=8$, respectively. Prior to being input to the DiT, this latent tensor is patchified using a kernel of $(1, 2, 2)$, and the resulting tokens serve as the input sequence. The umT5 model encodes the text prompt, and its embeddings are injected into the DiT via cross-attention to condition the generation. The model is trained to predict the velocity, which is used in the denoising step to recover the clean latent. Finally, the 3D VAE's decoder reconstructs this latent into the final video.

% \LHC{
\textbf{Wan2.1 model.} Wan2.1 is composed of three primary components: a 3D Variational Autoencoder (VAE), an umT5~\cite{chung2023unimax} text encoder, and a Diffusion Transformer (DiT). The 3D VAE compresses an input video of shape $(3, T, H, W)$ into a latent representation of shape $(16, T/t, H/h, W/w)$, where the temporal and spatial downsampling factors are $t=4$ and $h=w=8$, respectively. Prior to being input to the DiT, this latent tensor is patchified using a kernel of $(1, 2, 2)$, and the resulting tokens serve as the input sequence. The umT5 model encodes the text prompt, and its embeddings are injected into the DiT via cross-attention to condition the generation. The model is trained to predict the velocity, which is used in the denoising step to recover the clean latent. Finally, the decoder of 3D VAE reconstructs this latent into the final video.
% }

% Ace-step consists of a Music-DCAE (Deep Compression Autoencoder)~\cite{chen2024deep}, an umT5 text encoder, a lyric encoder, a speaker encoder, and a DiT. The raw audio waveform is first converted into a mel spectrogram. The Music-DCAE then encodes the input spectrogram of shape $(8, T, F)$ into a latent representation of shape $(8, T/t, F/f)$, with downsampling factors $t=f=8$. This latent is subsequently patchified using a kernel of $(16, 1)$ to produce the input tokens for the DiT. Conditional control is provided by three sources: the umT5 encoder for the music style prompt, the lyric encoder for the lyrics, and the speaker encoder for the speaker ID. These three embeddings are concatenated along the channel dimension and injected into the DiT via cross-attention. The model predicts the velocity to obtain the clean latent, which the Music-DCAE's decoder reconstructs into a mel spectrogram. A HiFiGAN vocoder~\cite{liao2024fish} is then used to convert this spectrogram into the final audio waveform.

% \LHC{
\textbf{Ace-step model.} Ace-step consists of a Music-DCAE (Deep Compression Autoencoder)~\cite{chen2024deep}, a umT5 text encoder, a lyric encoder, a speaker encoder, and a DiT. The raw audio waveform is first converted into a mel spectrogram. The Music-DCAE then encodes the input spectrogram of shape $(8, T, F)$ into a latent representation of shape $(8, T/t, F/f)$, with downsampling factors $t=f=8$ along the temporal and frequency axes. This latent is subsequently patchified using a kernel of $(16, 1)$ to produce the input tokens for the DiT. Conditional control is provided by three sources: the umT5 encoder for the music style prompt, the lyric encoder for the lyrics, and the speaker encoder for the speaker ID. These three embeddings are concatenated along the channel dimension and injected into the DiT via cross-attention. The model predicts the velocity to obtain the clean latent, which the Music-DCAE's decoder reconstructs into a mel spectrogram. A HiFiGAN vocoder~\cite{liao2024fish} is then used to convert this spectrogram into the final audio waveform.
% }

% The learning objective for both of the aforementioned models is \textit{Flow Matching}~\cite{lipman2022flow}. This framework trains a neural network, $v_\theta$, to predict a velocity field that transports samples from a simple source distribution, $p_0$ (e.g., Gaussian noise), to a complex target data distribution, $p_1$. Specifically, these models leverage Conditional Flow Matching. Given a noise sample $x_0 \sim p_0$ and a data sample $x_1 \sim p_1$, a simple linear interpolation path is defined for time $t$:
% $$
% x_t = (1-t)x_0 + t x_1
% $$
% The target velocity vector along this path is constant: $u_t = x_1 - x_0$. The model $v_\theta(x_t, t, c)$, conditioned on $c$, is trained to predict this vector by minimizing the following L2 loss:
% $$
% \mathcal{L}_{\text{FM}} = \mathbb{E}_{t \sim U(0,1), x_0 \sim p_0, x_1 \sim p_1} \left[ \| v_\theta((1-t)x_0 + t x_1, t, c) - (x_1 - x_0) \|^2 \right]
% $$
% This objective directly trains the model to learn the vector field that maps noise to data, which can lead to more stable and efficient training compared to traditional score-matching objectives.


% \LHC{
\textbf{Base model pre-training.} The learning objective for both of the aforementioned models is \textit{Flow Matching}~\cite{lipman2022flow}. This fashion trains a neural network, $v_\Theta(\cdot, t, c)$, to predict a velocity field that transports samples from a simple source distribution, $p_0$ (\textit{e.g.}, Gaussian noise), to a complex target data distribution, $p_1$. Specifically, these models leverage Conditional Flow Matching. Given a noise sample $x_0 \sim p_0$ and a data sample $x_1 \sim p_1$, a simple linear interpolation path is defined for time $t$:
\begin{equation}
    x_t = (1-t)x_0 + t x_1.
\end{equation}
The target velocity vector along this path is constant: $u_t = x_1 - x_0$. The model $v_\Theta(x_t, t, c)$, conditioned on $c$, is trained to predict this vector by minimizing the following L2 loss:
$$
\mathcal{L}_{\text{FM}} = \mathbb{E}_{t \sim U(0,1), x_0 \sim p_0, x_1 \sim p_1} \left[ \| v_\theta((1-t)x_0 + t x_1, t, c) - (x_1 - x_0) \|^2 \right].
$$
This objective directly trains the model to learn the vector field that maps noise to data, which leads to more stable and efficient training compared to traditional score-matching objectives.
% }

\subsection{Data Curation}
We curated a large-scale, high-quality dataset from a diverse range of sources to train our model. The primary component was sourced from YouTube, encompassing content such as music variety shows, classical music performances, cooking tutorials, public speeches, interviews, vlogs, and demonstrations of tool usage. This was supplemented with cinematic movie clips and high-quality stock footage from Pexels. To further bolster the audio modality, we also incorporated the widely-used VGGSound and AudioSet datasets.

For our self-collected data (YouTube, Pexels, and movie clips), we implemented a rigorous multi-stage filtering pipeline to ensure data quality and relevance:
\begin{itemize}[leftmargin=*]
    \item \textbf{Audio-Visual Pre-screening} Videos lacking an audio track were immediately discarded.
    \item \textbf{Quality Control} We filtered out content based on technical specifications: resolution below $1080p$, a bitrate-to-resolution ratio under $600$, and a DOVER~\cite{wu2023exploring} aesthetic quality score below $0.6$.
    \item \textbf{Temporal Coherence} PySceneDetect\footnote{https://github.com/Breakthrough/PySceneDetect} was applied to segment videos, and any resulting clip shorter than $5$ seconds was removed to ensure meaningful duration.
    \item \textbf{Audio Activity Detection} To eliminate silent segments, we analyzed each audio track for metrics such as volume, energy, and zero-crossing rate.
    \item \textbf{Speech Content Verification} Whisper~\cite{radford2023robust} was used to detect the presence of human speech. Clips without speech were retained as general audio-visual data. If speech was present, it proceeded to the next step.
    \item \textbf{Human Face Detection} For clips identified as containing speech, a second verification step was performed: we detected for the presence of a human face~\cite{deng2020retinaface}. If no face was found, the clip was discarded. If a face was present, we employed SyncNet~\cite{chung2016out} to verify the audio-visual correspondence (lip-sync). Only clips with a SyncNet confidence score above a threshold of $2.0$ were retained and explicitly labeled as containing speech content.
\end{itemize}


For the VGGSound and AudioSet data, a simplified process was used where we either performed scene detection or segmented clips based on their existing timestamps, retaining only those longer than $5$ seconds.

Following this comprehensive curation process, our final dataset comprises $7,685$ hours of data. This is categorized into three subsets: $1,187$ hours of verified speech-centric content, $3,074$ hours of general-purpose audio-video data, and $3,422$ hours from VGGSound and AudioSet primarily used for bolstering audio-specific training.

\subsection{Online Data Annotation}
Conventional offline annotation methods, where captions are pre-generated for entire videos, present a significant challenge for training generative models. During training, fixed-length clips are randomly sampled from these videos, often creating a temporal and semantic misalignment between the sampled clip and the global, pre-existing caption. This issue is particularly acute in the context of joint audio-video generation, where the temporal synchronization between an acoustic event and its description is critical. Even minor temporal shifts can render an audio annotation invalid.

To overcome these limitations, we propose and implement an Online Data Annotation Pipeline. This pipeline operates as a dedicated server process that runs concurrently with training. It dynamically fetches raw video, processes clips in real-time to generate precisely aligned data-annotation pairs, and populates a shared buffer. The main training process then acts as a consumer, fetching these ready-to-use data tuples, ensuring that every training instance is perfectly synchronized.

The online processing for each data tuple involves the following steps:
\begin{itemize}[leftmargin=*]
    \item \textbf{Temporal Sampling} A fixed-length segment (e.g., $5$ seconds) is randomly extracted from a source video, yielding corresponding video and audio streams.
    \item \textbf{Multi-modal Annotation} The extracted audio-video clip is immediately passed to our annotation module for captioning.
    \item \textbf{Text and Video Encoding} The generated text prompts (for video, audio, and speech) are encoded. Concurrently, the video clip is encoded into a spatio-temporal latent representation using the 3D VAE.
    \item \textbf{Audio Encoding} The audio stream is converted to a mel spectrogram and subsequently encoded into a latent representation by the Music-DCAE~\cite{chen2024deep}.
\end{itemize}

The core of our pipeline is the multi-modal annotation step (Step $2$), which proceeds as follows:
\begin{itemize}[leftmargin=*]
    \item \textbf{Speech Transcription} Whisper~\cite{radford2023robust} is employed to perform Automatic Speech Recognition (ASR) on the sampled audio, yielding the raw speech content.
    \item \textbf{Structured Multimodal Captioning} We construct a structured prompt that incorporates the transcribed speech. This prompt, along with the audio and video streams of the clip, is fed into the QWen2.5-Omni~\cite{xu2025qwen2} multimodal model. QWen2.5-Omni is specifically instructed to output three distinct, aligned annotations for the clip: the verified speech content, a descriptive video caption, and a caption for the ambient audio.
\end{itemize}

This online, just-in-time process guarantees that every training instance consists of video and audio latents that are perfectly synchronized in time and semantically consistent with their corresponding textual annotations, thereby eliminating the data misalignment problem inherent in offline methods.

\begin{figure}[t!]
  % \vspace{-1em}
  \centering
  \includegraphics[width=\linewidth]{figures/universe_final_version4_compressed.pdf}
  \caption{\textbf{Architecture of UniVerse-1.} (a) Overall architecture. The architectural foundation of UniVerse-1 is realized through a stitching of experts methodology. This approach deeply integrates the pre-trained Wan2.1 video model and the Ace-step audio model. (b) Fused block. The fusion is implemented at a granular, block-by-block level, where each block in the Wan architecture is deeply fused with its corresponding block in the Ace-step architecture.}
  \label{fig:archi}
\end{figure}


\begin{figure}[t!]
  % \vspace{-1em}
  \centering
  \includegraphics[width=\linewidth]{figures/universe_attn.pdf}
  \caption{\textbf{Revised attention of UniVerse-1.} (a) Self attention of video branch, with additional mel tokens as input. (b) Cross attention of video branch, with additional mel tokens as input. (c) Cross attention of mel branch, with addition video tokens as input.}
  \label{fig:attn}
\end{figure}

\subsection{Method}
The overall architecture of our method is depicted in Fig.~\ref{fig:archi}. We introduce several targeted modifications to the input stages of the original Wan2.1 and Ace-step models to facilitate bimodal integration and control.
For the video component (Wan2.1), we enable conditioning on a reference image. During the forward process, the first frame of the noisy video latent is replaced with the corresponding clean latent representation of the provided reference image.
For the audio component (Ace-step), we perform two adjustments. First, to ensure temporal alignment with the video's 25 frames-per-second (fps) rate, the input mel spectrograms are processed at a 25.6 kHz sampling rate instead of the original 44.1 kHz. Second, to generalize the model beyond speaker-specific generation, we have removed the speaker encoder and its corresponding input from the architecture.

\subsubsection{Stitching of experts} We introduce a novel framework, termed ``Stitching of experts'', for integrating specialized, pre-existing models for video and audio synthesis. This approach is designed to preserve the generative capabilities of each unimodal expert while simultaneously enabling fine-grained, bidirectional interaction between them at the level of individual layer blocks. To enhance training efficiency and leverage the powerful priors of these pre-trained models, we apply the stitching technique to the Wan2.1 and Ace-step models at the transformer block level. This process results in a unified, dual-stream architecture where each block co-processes information from both the video and audio modalities, functioning akin to a Mixture-of-Experts (MoE) layer. \\
As illustrated in Fig~\ref{fig:archi}. b), we facilitate bidirectional cross-modal communication within each block. Specifically, the hidden states from the video stream, following its self-attention module, are injected into the audio stream's cross-attention module. Conversely, the hidden states from the audio stream, following its linear attention module, are reciprocally injected into the video stream's self-attention and cross-attention module. To ensure consistent feature scaling, the hidden states from both streams are jointly passed through a shared LayerNorm layer before injected into video stream's attention.\\
Prior to injection, the cross-modal hidden states are passed through a two-layer linear adapter for feature space alignment. In the video stream, for instance, features from the audio branch are projected using dedicated key ($k_{proj}$) and value ($v_{proj}$) layers~\ref{fig:attn}(a). A frame-by-frame cross-attention is then performed with the queries from the video stream to ensure alignment between the video and audio. Within each stream's respective cross-attention mechanism(shown in ~\ref{fig:attn}(b) and ~\ref{fig:attn}(c)), the conditioning signal from the other modality is projected using dedicated key ($k_{proj}$) and value ($v_{proj}$) layers. These new key-value pairs are then concatenated with the original text-derived key-value pairs along the context dimension, thereby enriching the conditioning information with cross-modal context.
\subsubsection{Layer Interpolation} A key challenge in stitching the Wan2.1 and Ace-step models is the architectural mismatch in their depth, as they possess a different number of transformer blocks. To reconcile this disparity, we introduce a \textbf{layer interpolation technique}.
This method involves first calculating the difference in the number of layers. We then strategically insert new blocks at uniform intervals into the shallower of the two models until their depths align. Crucially, the parameters for each new block are initialized by linearly interpolating the weights of its immediately adjacent (bracketing) layers. This initialization strategy effectively bridges the architectural gap while ensuring a smooth performance trajectory during training, thereby mitigating the risk of training instability and severe performance oscillations.
\subsubsection{Training Loss}
In addition to the primary Flow Matching objective (Sec.~\ref{sec:preliminary}), we incorporate two additional loss functions.
\paragraph{Semantic Alignment Loss}
For the audio modality, we employ a \textbf{Semantic Similarity Loss} ($\mathcal{L}_{\text{SSL}}$), a technique consistent with Ace-step, to enhance the semantic fidelity of the generated audio. This loss operates by aligning an intermediate feature representation, $h_{\text{audio}}$, extracted from the audio stream of our fused block at a specific layer $L$ ($L=12$ in our configuration). This internal representation is aligned against target features derived from two expert, pre-trained audio models:

\begin{itemize}[leftmargin=*]
    \item \textbf{MERT}(Music Encoder Representations from Transformers)~\cite{li2023mert}, which provides a general musical representation, $h_{\text{mert}}$, with a dimensionality of 1024 x $T_m$ (at a 75 Hz frame rate).
    \item \textbf{mHuBERT}(multilingual HuBERT)~\cite{boito2024mhubert}, which provides a speech-centric representation, $h_{\text{mHuBERT}}$, with a dimensionality of 768 x $T_h$ (at a 50 Hz frame rate).
\end{itemize}

To compute this loss, the intermediate feature $h_{\text{audio}}$ is first processed by two separate projection heads ($\pi_{\text{MERT}}$ and $\pi_{\text{mHuBERT}}$) and temporally interpolated to match the dimensionality and sequence length of $h_{\text{MERT}}$ and $h_{\text{mHuBERT}}$, respectively. The semantic similarity loss is then defined as the negative cosine similarity, encouraging the model's internal representations to align with those of the expert models:
$$
\mathcal{L}_{\text{SSL}} = \frac{1}{2} (\text{cosineSim}(h'_{\text{audio}}, h'_{\text{MERT}}) + \text{cosineSim}(h'_{\text{audio}}, h'_{\text{mHuBERT}}))
$$
where $h'$ denotes the temporally aligned representations.

\paragraph{Low Quality Data Loss Strategy} 

The AudioSet and VGGSound datasets, while offering rich auditory diversity, are characterized by low visual fidelity. To leverage their strong audio content without corrupting the video generation quality, we employ a conditional loss scheme.
Specifically, the Flow Matching loss for the video modality ($\mathcal{L}_{\text{FM-video}}$) is only computed for samples originating from these two datasets when the diffusion timestep \textit{t} exceeds an empirically determined threshold. We set this threshold to $\tau = 800$ (out of 1000 total timesteps). This strategy is predicated on the principle that at high noise levels (i.e., for $t > \tau$), the model learns to capture coarse, low-frequency features of the video, such as general motion and structure, which are less affected by the poor visual quality. By excluding the loss calculation at lower noise levels, we prevent the model from overfitting to the high-frequency visual artifacts and noise present in these datasets. The loss for video modality is:
$$
\mathcal{L}_{\text{FM-video}} = \left\{
\begin{array}{rcl}
\mathcal{L}_{\text{FM}}(x_1, x_0, t) & & if ~~ x_1 \in \Theta ~~ or ~~ (x_1 \in \zeta ~~ and ~~ t > 800)\\
0~~~~~~~~~ & & else
\end{array} \right.
$$
where $x_0 \sim p_0$ is noise sample, $x_1 \sim p_1$ is data sample, $\zeta$ is data subset include vggSound and audioset, $\Theta$ is data subset exclude vggSound and audioset. The final training objective is a weighted sum of the flow matching and semantic alignment losses:
$$
\mathcal{L} = \mathcal{L}_{\text{FM-video}} + \mathcal{L}_{\text{FM-mel}} + \lambda_{SSL}\cdot\mathcal{L}_{\text{SSL}}
$$
where $\mathcal{L}_{\text{SSL}}$ is a hyperparameter controlling the influence of the SSL guidance, empirically set to $1.0$ according to Ace-step.

\subsection{Independent Noise Sampling Strategy}

Our empirical investigation reveals a critical sensitivity of multi-modal diffusion models to the pseudo-random number generation process. When a single, fixed random seed is used to initialize a training run, the noise tensors for the video (${\epsilon}_v$) and audio (${\epsilon}_a$) modalities are sampled sequentially from the same deterministic PRNG sequence. Due to the deterministic nature of the underlying algorithm (Linear Congruential method), this sequential sampling introduces a spurious structural correlation between the two noise tensors, which can be expressed as ${\epsilon}_a = f({\epsilon}_v)$. This violates the critical assumption that the noise vectors are statistically independent.

The model inadvertently learns this spurious correlation as a shortcut during training. Consequently, during inference, any alteration to the sampling of ${\epsilon}_v$, such as a change in video resolution or duration, propagates through the PRNG's state and alters the structure of the subsequently sampled ${\epsilon}_a$. This mismatch with the learned correlation results in a significant degradation of the audio generation quality.

To address this, we propose an \textbf{Independent Noise Sampling Strategy}. This approach isolates the noise generation for each modality by employing separate and independently seeded PRNG instances. This method effectively breaks the deterministic correlation, ensuring the noise vectors are statistically independent. As a result, the model becomes robust to variations in inference-time conditions, mitigating the issue of performance degradation.


\section{Experiments}
\subsection{Setup}
\paragraph{Implementation Details}
The training is conducted with an effective batch size of $128$ over $50k$ steps on a $7,600$-hour audio-visual datasets built by our data curation pipeline,
using the AdamW optimizer with a learning rate of $5e-6$. We employ Fully Sharded Data Parallel (FSDP) for distributed training across multiple nodes, with a gradient accumulation step of $4$.
\paragraph{Compared Methods} We conduct a comprehensive evaluation of our model by benchmarking it against a suite of state-of-the-art baselines across several distinct generation tasks:
\begin{itemize}[leftmargin=*]
    \item Video Generation: We compare our model against leading text-to-video systems, including Wan2.2 (14B)~\cite{wan2025wan}, HunyuanVideo~\cite{kong2024hunyuanvideo}, CogVideoX-1.5, (5B)~\cite{yang2024cogvideox}, Kling 2.1~\cite{kuaishou2024kling}, and SeeDance 1.0~\cite{gao2025seedance}.
    \item Audio Generation: For the audio modality, we benchmark against established text-to-audio models such as Stable Audio Open~\cite{evans2025stable} and AudioLDM2~\cite{audioldm2-2024taslp}.
    \item Text-to-Speech (TTS): As a supplementary evaluation of vocal synthesis, we compare our model's performance against specialized TTS models, including CosyVoice~\cite{du2024cosyvoice}, CosyVoice2~\cite{du2024cosyvoice2}, and VibeVoice~\cite{peng2025vibevoice}.
    \item Audio-to-Video methods: We also compare out model with talking-based audio to video methods such as FantastyTalking~\cite{wang2025fantasytalking} and Wan-S2V~\cite{gao2025wan}.
    \item Video-to-Audio methods: As a closely related task to joint audio-visual generation, we also benchmarked our model on the video-to-audio (V2A) task, drawing comparisons with state-of-the-art methods such as HunyuanVideo-Foley~\cite{shan2025hunyuanvideo}.
    \item Joint Audio-Video Generation: For our core task of synchronous audio-visual synthesis, we compare our method against existing prompt-based joint generation models: SVG~\cite{ishii2024simple}. Since methods such as MM-Diffusion~\cite{ruan2023mm} and R-FLAV~\cite{ergasti2025r} are class-conditional generative models, a direct comparison with our approach is not applicable, we only compare with SVG in our report.
\end{itemize}
It is important to note that our model is constructed by stitching the pre-trained WAN2.1 (1.3B) and Step-Ace (3.5B) models. Consequently, the comparisons against the aforementioned state-of-the-art baselines are intended to provide a qualitative reference and situate our work, rather than to make a direct claim of superior performance.

\paragraph{Benchmark} To construct our evaluation set, we curated $600$ image-text prompt pairs from a multitude of sources. These sources encompass frames extracted from YouTube videos, BiliBili videos, TikTok clips, movies, and anime; images generated by AI models~\cite{jimeng,wu2025qwen}; and a collection of images from public websites. Our dataset comprises three subsets.
\begin{itemize}[leftmargin=*]
    \item Set1-I contains image-text pairs (including AI-generated, web-crawled, and media screenshots), for which video/audio captions and speech content were produced using LLMs~\cite{xu2025qwen2} and manual annotation, comprising a total of $205$ samples. Statistical results in~\ref{fig:unibench_statistics_set1}.
    \item Set2-V consists of video clips from YouTube and Bilibili, which were annotated with LLM-generated~\cite{xu2025qwen2} captions and Whisper-based ASR~\cite{radford2023robust} transcripts, followed by human verification, comprising a total of $295$ samples. Statistical results in~\ref{fig:unibench_statistics_set2}.
    \item Set3-Ted (the Verse-Ted subset) includes TED Talks from September 2025, processed with the same annotation pipeline as Set2, comprising a total of $100$ samples. 
\end{itemize}
Our collected test data is highly diverse, encompassing a wide spectrum of audio categories: human speech; animal vocalizations (e.g., bird chirping, cat meowing); instrumental music (e.g., piano, guitar); natural sounds (e.g., thunder, rain); human-object interactions (e.g., keyboard typing, cooking, chopping vegetables); object-object interactions (e.g., glass shattering, marbles dropping); mechanical sounds (e.g., trains, airplanes), and so on. The complete statistical results of set1 and set2 are shown in~\ref{fig:unibench_statistics_all}.

\begin{figure*}
    \centering
    \begin{subfigure}[h]{0.3\textwidth}
    \centering
    \includegraphics[width=1.0\linewidth]{figures/audio_statistics_all.pdf}
    % \vspace{-3mm}
    \caption{Statistical results of set1/2.}
    \label{fig:unibench_statistics_all}
    % \vspace{-5mm}
    \end{subfigure}
    \hspace{0.01\textwidth} % 可选的右侧间距
    \vrule
    \hspace{0.01\textwidth} % 可选的左侧间距
    \begin{subfigure}[h]{0.3\textwidth}
    \centering
    \includegraphics[width=1.0\linewidth]{figures/audio_statistics_set1.pdf}
    % \vspace{-3mm}
    \caption{Statistical results of set1.}
    \label{fig:unibench_statistics_set1}
    % \vspace{-5mm}
    \end{subfigure}
    \hspace{0.01\textwidth} % 可选的右侧间距
    \vrule
    \hspace{0.01\textwidth} % 可选的左侧间距
    \begin{subfigure}[h]{0.3\textwidth}
    \centering
    \includegraphics[width=1.0\textwidth]{figures/audio_statistics_set2.pdf}
    % \vspace{-3mm}
    \caption{Statistical results of set2.}
    \label{fig:unibench_statistics_set2}
    \end{subfigure}
    
    \caption{Statistical results of \textbf{Verse-Bench}. Best viewed with \textbf{zoom-in}. A larger, high-resolution version is available in Appendix~\ref{lab:verse-bench-fig}}
    % \vspace{-7mm}
\end{figure*}

\paragraph{Evaluation Protocol}
We quantitatively evaluate our method against baselines on the Verse-Bench benchmark. The evaluation is structured across $6$ distinct generation tasks, each with a tailored set of metrics.

\begin{itemize}[leftmargin=*]
    \item Video Generation: Performance is assessed on three criteria:
    \begin{itemize}
        \item Motion Score (MS): This metric quantifies motion dynamics, calculated from the normalized optical flow magnitude detected by the RAFT model~\cite{teed2020raft}.
        \item Aesthetic Score (AS): This is a composite score averaging three components: fidelity, measured by MANIQA~\cite{yang2022maniqa} to penalize blur and artifacts, and aesthetic quality, evaluated by both aesthetic-predictor-v2-5~\cite{aesthetic-predictor-v2-5} and Musiq~\cite{ke2021musiq}.
        \item ID Consistency (ID): To measure identity preservation, we compute the mean DINOV3~\cite{simeoni2025dinov3} feature similarity between the reference image and each generated frame.
    \end{itemize} 
        
    \item Audio Generation: We evaluate audio quality from three perspectives:
    \begin{itemize}
        \item Distributional Similarity: We measure the Fréchet Distance (FD) and Kullback-Leibler (KL) divergence between the generated and real data distributions, using features extracted from PANNs~\cite{kong2020panns} and PaSST~\cite{koutini2021efficient}.
        \item Semantic Consistency: The alignment between the audio and the input text is measured by the LAION-CLAP~\cite{wu2023large} score.
        \item Quality and Diversity: We report the Inception Score (IS) calculated with a PANNs classifier. Additionally, we use AudioBox-Aesthetics~\cite{tjandra2025meta} to assess Production Quality (PQ), Production Complexity (PC), Content Enjoyment (CE), and Content Usefulness (CU).
    \end{itemize}
    \item Text-to-Speech (TTS): We evaluate synthesis accuracy using the Word Error Rate (WER), which is derived by transcribing the generated audio with the Whisper-large-v3 model~\cite{radford2023robust}.
    \item Audio-to-Video: We evaluate this task using the same criteria as the video generation task, additionally providing a SyncNet~\cite{chung2016out} confidence score to assess lip-sync accuracy.
    \item Video-to-Audio: This task use all metrics from audio generation tasks. Furthermore, we introduce the Audio-Video Alignment (AV-A) metric to specifically quantify the temporal synchronization between the generated audio and video streams, which is computed via Synchformer~\cite{iashin2024synchformer}.
    \item Joint Audio-Video Generation: For this task, we use all relevant metrics from the individual tasks above. 
\end{itemize}

The evaluation of our models and their components is conducted across the three test sets as follows:
\begin{itemize}[leftmargin=*]
    \item The video generation models and SVG are evaluated on Set 1 and Set 2.
    \item For the audio generation model is evaluated on Set1 and Set2.
    \item The Text-to-Speech (TTS) model is primarily evaluated on Set 3.
    \item The Audio-to-Video (A2V) model is evaluated exclusively on Set 3, while the Video-to-Audio (V2A) model is evaluated exclusively on Set 2.
    \item Finally, our complete Universe-1 model is benchmarked against all three test sets.
\end{itemize}
For audio generation, we evaluate the metrics CE, CU, PC, and PQ on Set 1. On Set 2, the evaluation is expanded to include FD, KL, and CS in addition to the aforementioned metrics. Furthermore, LSE-C is evaluated exclusively on Set 3, while the AV-A metric is also applied to Set 1 when evaluating UniVerse-1 and SVG.

\subsection{Quantitative Evaluation}
% We compared Universe-1 with various state-of-the-art (SOTA) models as shown in Tab.~\ref{tab:uni_bench}. In video generation, our model surpasses all competing methods in identity preservation (ID: 0.89) and achieves a competitive score in action similarity (AS: 0.47). For audio generation, while not outperforming specialized audio models on all metrics, it shows strong performance in pitch correlation (PC: 2.49), surpassing most baselines. The most significant advantage of our model is demonstrated in audio-video coherence tasks, where its lip-sync score (LSE-C: 7.10) is substantially higher than other methods, strongly validating the effectiveness of our joint generation framework. Finally, in the TTS task, the Word Error Rate (WER) of 0.18 indicates performance comparable to that of dedicated TTS models.our methods. 

We compare Universe-1 against a range of state-of-the-art (SOTA) models specialized for specific tasks as shown in Tab.~\ref{tab:uni_bench}. As a unified joint generation model, a direct comparison of single-modality metrics against these "expert" models presents inherent complexities. Nevertheless, our model demonstrates robust capabilities across multiple dimensions.
In terms of video quality, Universe-1 achieves the highest score in identity preservation (ID: 0.89), showcasing its superior ability to maintain subject consistency throughout the generation process. For audio quality, while there is a gap compared to leading audio-only generation models, our model obtains a highly competitive score in pitch correlation (PC: 2.49).
The core strength of our model lies in synchronous audio-video generation. It is crucial to interpret the metrics for such joint generation tasks with caution. For instance, in the video-to-audio (V2A) setting, the AV-A metric for a model like Hunyuanvideo-Foley is calculated with a ground-truth video, whereas our model generates both modalities simultaneously. Therefore, AV-A must be considered in conjunction with the audio-text CLAP score (CS) for a holistic assessment. From this integrated perspective, our model (AV-A: 0.23, CS: 0.16) demonstrates a better overall audio-visual content consistency than SVG (AV-A: 0.09, CS: 0.08).
Similarly, for the lip-sync (LSE-C) metric, our model's score (1.34) is evaluated on fully generated audio and video, making it susceptible to the quality of both generated modalities. In contrast, audio-to-video (A2V) methods like Wan-S2V are evaluated using ground-truth audio, which naturally yields a higher score (6.49). Despite this evaluation disparity, Universe-1 achieves promising results as the first open-source joint generation framework of its kind, establishing a solid foundation for future research.

\begin{table}
\vspace{-1em}
\small
\centering
\resizebox{\linewidth}{!}{
    \renewcommand{\arraystretch}{1.8}
    \begin{tabular}{l|c|c|c|c|c|c|c|c|c|c|c|c|c|c}
    \toprule
    % --- 表头部分 ---
    % 第一行主分类表头，使用 \multicolumn 合并单元格
    \multicolumn{2}{c|}{} & \multicolumn{3}{c|}{video} & \multicolumn{7}{c|}{audio} & \multicolumn{1}{c|}{tts} & \multicolumn{2}{c}{Audio-Video} \\
    \midrule
    % 第二行子分类表头
     & method & MS$~\uparrow$ & AS$~\uparrow$ & ID$~\uparrow$ & FD$~\downarrow$ & KL$~\downarrow$ & CS$~\uparrow$ & CE$~\uparrow$ & CU$~\uparrow$ & PC$~\downarrow$ & PQ$~\uparrow$ & WER$~\downarrow$ & LSE-C$~\uparrow$ & AV-A$~\downarrow$ \\
    \midrule
    
    % --- 数据行部分 ---
    % Video 部分，跨越5行
    \multirow{5}{*}{\makecell{video \\ methods}} 
     & CogVideox1.5 5B~\cite{yang2024cogvideox}         & 0.43 & 0.44 & 0.83 & - & - & - & - & - & - & - & - & - & - \\ \cline{2-15} 
     & Wan2.2-14B~\cite{wan2025wan}                     & \textbf{1.04} & \textbf{0.50} & \underline{0.88} & - & - & - & - & - & - & - & - & - & - \\ \cline{2-15} 
     & Kling2.1~\cite{kuaishou2024kling}                & 0.31 & 0.41 & 0.85 & - & - & - & - & - & - & - & - & - & - \\ \cline{2-15}
     & SeeDance1.0~\cite{gao2025seedance}               & \underline{0.50} & \underline{0.47} & 0.86 & - & - & - & - & - & - & - & - & - & - \\
    \midrule
    % Audio 部分，跨越2行
    \multirow{2}{*}{\makecell{audio \\ methods}} 
     & Stable Audio~\cite{evans2025stable}              & - & - & - & \underline{1.13} & \underline{1.38} & \underline{0.28} & 3.88 & \textbf{6.15} & \underline{2.60} & 6.53 & - & - & - \\ \cline{2-15} 
     & AudioLdm2~\cite{audioldm2-2024taslp}             & - & - & - & 1.21 & 2.30 & 0.24 & \underline{3.98} & \underline{5.88} & 3.43 & 6.04 & - & - & - \\
    \midrule
    % TTS 部分，跨越3行
    \multirow{3}{*}{\makecell{tts \\ methods}}   
     & Cosyvoice~\cite{du2024cosyvoice}                 & - & - & - & - & - & - & - & - & - & - & 0.17 & - & - \\ \cline{2-15}  
     & Cosyvoice2~\cite{du2024cosyvoice2}               & - & - & - & - & - & - & - & - & - & - & \underline{0.16} & - & - \\ \cline{2-15}
     & Vibevoice~\cite{peng2025vibevoice}               & - & - & - & - & - & - & - & - & - & - & \textbf{0.15} & - & - \\
    \midrule
    % Audio2Video 部分，跨越2行
    \multirow{2}{*}{\makecell{a2v \\ methods}}   
     & FantastyTalking~\cite{wang2025fantasytalking}    & 0.07 & 0.42 & 0.87 & - & - & - & - & - & - & - & - & 2.68 & - \\ \cline{2-15}  
     & Wan-S2V~\cite{gao2025wan}                        & 0.17 & 0.46 & \textbf{0.89} & - & - & - & - & - & - & - & - & 6.49 & - \\
    \midrule
    % Video2Audio 部分，跨越1行
    \multirow{1}{*}{\makecell{v2a \\ methods}}     
     & Hunyuanvideo-Foley~\cite{shan2025hunyuanvideo}   & - & - & - & \textbf{0.82} & \textbf{1.27} & \textbf{0.40} & \textbf{4.04} & 5.72 & 3.09 & 6.27 & - & - & 0.78 \\
    \midrule
    % AV joint 部分，跨越2行
    \multirow{2}{*}{\makecell{joint \\ methods}}  
     & SVG~\cite{ishii2024simple}                       & 0.40 & 0.41 & 0.25 & 1.55 & 3.62 & 0.08 & 2.93 & 5.50 & \textbf{2.35} & 6.26 & - & - & 0.09 \\ \cline{2-15} 
     & \colorbox{CornflowerBlue}{\textbf{Ours}}         & 0.20 & 0.47 & 0.89 & 1.25 & 2.70 & 0.16 & 3.53 & 4.61 & 2.49 & 5.20 & 0.18 & 1.34 & 0.23 \\ 
    \bottomrule
    \end{tabular}
}
\vspace{1em}
\caption{Quantitative results compared with baselines on Verse-Bench.}
\label{tab:uni_bench}
\vspace{-1em}
\end{table}


% \subsection{Qualitative Evaluation}
% \subsection{User Study}
\subsection{Ablation Study}
We performed an ablation study to investigate the contributions of our Low Quality data Loss Strategy(LQLS) and Independent Noise Sampling Strategy(INSS), as shown in Tab.~\ref{tab:ablation}. The findings indicate that LQLS provides improvement accross video quality and consistency ID, thus confirming its efficacy and positive impact on training. Furthermore, the results for INSS demonstrate a significant enhancement in audio generation quality, validating the effectiveness of this approach.

\begin{table}
% \vspace{-1em}
\small
\centering
\resizebox{\linewidth}{!}{
    \renewcommand{\arraystretch}{1.8}
    \begin{tabular}{l|c|c|c|c|c|c|c|c|c|c|c|c|c}
    \toprule
    % --- 表头部分 ---
    method & MS$~\uparrow$ & AS$~\uparrow$ & ID$~\uparrow$ & FD$~\downarrow$ & KL$~\downarrow$ & CS$~\uparrow$ & CE$~\uparrow$ & CU$~\uparrow$ & PC$~\downarrow$ & PQ$~\uparrow$ & WER$~\downarrow$ & LSE-C$~\uparrow$ & AV-A$~\downarrow$ \\
    \midrule
    % --- 数据行部分 ---
     w/o LQLS                                  & 0.38 & 0.44 & 0.78 & 1.26 & 2.84 & 0.15 & 3.38 & 4.35 & 2.58 & 4.97 & 0.16 & 1.35 & 0.28 \\ \cline{1-14}
     w/o INSS                                  & 1.10 & 0.43 & 0.75 & 1.43 & 3.51 & 0.11 & 2.44 & 3.14 & 2.92 & 3.99 & 0.38 & 0.99 & 0.18 \\ \cline{1-14}
     \colorbox{CornflowerBlue}{\textbf{Ours}}  & 0.20 & 0.47 & 0.89 & 1.25 & 2.70 & 0.16 & 3.53 & 4.61 & 2.49 & 5.20 & 0.18 & 1.34 & 0.23 \\
    \bottomrule
    \end{tabular}
}
\vspace{1em}
\caption{Ablation results on Verse-Bench.}
\label{tab:ablation}
\vspace{-1em}
\end{table}
\section{Limitation and Future Work}
The work presented herein constitutes an initial exploration into unified audio-video generation. Our study was constrained by computational resources, necessitating that we conduct training exclusively on the Wan2.1-1.3B video model. The performance of our model is, therefore, inherently limited by the capacity of this base model.
In the future, our research will focus on two key directions. First, we will scale up our experiments to larger video foundation models. Second, we will engage in more extensive and refined data curation efforts. The ultimate objective is to significantly advance the capabilities of open-source audio-video synthesis models, thereby bridging the performance gap to state-of-the-art proprietary models.
\section{Conclusion}

In this paper, we presented UniVerse-1, a novel framework for joint audio-video synthesis, achieved through the deep integration of a video foundation model and a music generation model using stitching of experts. Following fine-tuning on the dataset we curated, we also introduced Verse-Bench, a comprehensive benchmark to foster comparative research. To promote reproducibility and further innovation, our model and code have been made publicly available. Our experimental results validate that this methodology offers a viable and efficient pathway for building sophisticated multimodal generative models by leveraging pre-existing unimodal foundations.

\newpage

{
    \small
    \bibliographystyle{unsrt}
    \bibliography{neurips_2025}
}

\newpage
\appendix
\section*{Appendix}
\addcontentsline{toc}{section}{Appendix}
\section{More Details about Verse-Bench}
\subsection{Statistical Results}
\label{lab:verse-bench-fig}
This section provides the detailed category catalog for Verse-Bench, along with a high-resolution pie chart illustrating its statistical distribution in Fig.~\ref{fig:appendix_unibench_statistics_all}~\ref{fig:appendix_unibench_statistics_set1}~\ref{fig:appendix_unibench_statistics_set2}.

Category list is in Tab.~\ref{tab:audio_stats_set1},~\ref{tab:audio_stats_set2},~\ref{tab:audio_stats_combined}.

%%%%%%%%%%%%%%%%%%%%%%%%%%%%%%%%%%%%%%%%%%%%%%%%%%%%
\begin{table}[h!]
    \centering
    \caption{Detailed Audio Classification Statistics for Set 1}
    \label{tab:audio_stats_set1}
    \begin{tabularx}{\linewidth}{l S[table-format=2.0] S[table-format=2.1]}
        \toprule
        \textbf{Category / Subcategory} & \textbf{Count} & {\textbf{Percentage (\%)}} \\
        \midrule
        
        \textbf{Natural Environment} & \textbf{70} & \textbf{36.3} \\
        \quad Weather-Wind & 18 & {} \\
        \quad Plants-Vegetation & 12 & {} \\
        \quad Animals-Birds & 10 & {} \\
        \quad Water-Ocean/Waves & 8 & {} \\
        \quad \textit{Other Natural Subcategories} & 22 & {} \\ % Combines smaller subcategories
        
        \addlinespace % Adds a little vertical space between major categories
        
        \textbf{Music \& Instruments} & \textbf{40} & \textbf{20.7} \\
        \quad Musical Compositions & 14 & {} \\
        \quad Keyboard-Piano & 7 & {} \\
        \quad String-Guitar & 7 & {} \\
        \quad \textit{Other Music Subcategories} & 12 & {} \\
        
        \addlinespace
        
        \textbf{Daily Life} & \textbf{22} & \textbf{11.4} \\
        \quad Office-Writing/Typing & 4 & {} \\
        \quad Tools-Electronics & 4 & {} \\
        \quad \textit{Other Daily Life Subcategories} & 14 & {} \\
        
        \addlinespace
        
        \textbf{Human Voices} & \textbf{21} & \textbf{10.9} \\
        \quad Adults-Conversation & 6 & {} \\
        \quad \textit{Other Human Voices Subcategories} & 15 & {} \\
        
        \addlinespace
        
        \textbf{Transportation} & \textbf{12} & \textbf{6.2} \\
        
        \addlinespace
        
        \textbf{Industrial \& Urban} & \textbf{10} & \textbf{5.2} \\
        
        \addlinespace
        
        \textbf{Special Effects} & \textbf{9} & \textbf{4.7} \\
        
        \addlinespace
        
        \textbf{Weapons \& Explosions} & \textbf{7} & \textbf{3.6} \\

        \midrule
        
        \textbf{Total Classified} & \textbf{191} & \textbf{93.2} \\
        Unclassified & 14 & 6.8 \\
        \midrule
        \textbf{Grand Total} & \textbf{205} & \textbf{100.0} \\
        
        \bottomrule
    \end{tabularx}
\end{table}
%%%%%%%%%%%%%%%%%%%%%%%%%%%%%%%%%%%%%%%%%%%%%%%%%%%%

%%%%%%%%%%%%%%%%%%%%%%%%%%%%%%%%%%%%%%%%%%%%%%%%%%%%
\begin{table}[h!]
    \centering
    \caption{Detailed Audio Classification Statistics for Set 2}
    \label{tab:audio_stats_set2}
    \begin{tabularx}{\linewidth}{l S[table-format=3.0] S[table-format=2.1]}
        \toprule
        \textbf{Category / Subcategory} & \textbf{Count} & {\textbf{Percentage (\%)}} \\
        \midrule
        
        \textbf{Natural Environment} & \textbf{106} & \textbf{35.9} \\
        \quad Weather-Wind & 22 & {} \\
        \quad Animals-Birds & 22 & {} \\
        \quad Weather-Rain & 19 & {} \\
        \quad Animals-Dogs & 9 & {} \\
        \quad \textit{Other Natural Subcategories} & 34 & {} \\
        
        \addlinespace
        
        \textbf{Music \& Instruments} & \textbf{54} & \textbf{18.3} \\
        \quad Musical Compositions & 16 & {} \\
        \quad Keyboard-Piano & 10 & {} \\
        \quad String-Violin & 8 & {} \\
        \quad \textit{Other Music Subcategories} & 20 & {} \\
        
        \addlinespace
        
        \textbf{Daily Life} & \textbf{28} & \textbf{9.5} \\
        \quad Office-Writing/Typing & 6 & {} \\
        \quad Tools-Power Tools & 6 & {} \\
        \quad Communication-Phone & 6 & {} \\
        \quad \textit{Other Daily Life Subcategories} & 10 & {} \\
        
        \addlinespace
        
        \textbf{Human Voices} & \textbf{16} & \textbf{5.4} \\
        \quad Adults-Conversation & 9 & {} \\
        \quad \textit{Other Human Voices Subcategories} & 7 & {} \\
        
        \addlinespace
        
        \textbf{Transportation} & \textbf{10} & \textbf{3.4} \\
        
        \addlinespace
        
        \textbf{Industrial \& Urban} & \textbf{9} & \textbf{3.1} \\
        
        \addlinespace
        
        \textbf{Weapons \& Explosions} & \textbf{5} & \textbf{1.7} \\

        \addlinespace

        \textbf{Special Effects} & \textbf{2} & \textbf{0.7} \\

        \midrule
        
        \textbf{Total Classified} & \textbf{230} & \textbf{78.0} \\
        Unclassified & 65 & 22.0 \\
        \midrule
        \textbf{Grand Total} & \textbf{295} & \textbf{100.0} \\
        
        \bottomrule
    \end{tabularx}
\end{table}
%%%%%%%%%%%%%%%%%%%%%%%%%%%%%%%%%%%%%%%%%%%%%%%%%%%%

%%%%%%%%%%%%%%%%%%%%%%%%%%%%%%%%%%%%%%%%%%%%%%%%%%%%
\begin{table}[h!]
    \centering
    \caption{Combined Audio Classification Statistics for Set 1 \& Set 2}
    \label{tab:audio_stats_combined}
    \begin{tabularx}{\linewidth}{l S[table-format=3.0] S[table-format=2.1]}
        \toprule
        \textbf{Category / Subcategory} & \textbf{Count} & {\textbf{Percentage (\%)}} \\
        \midrule
        
        \textbf{Natural Environment} & \textbf{176} & \textbf{36.1} \\
        \quad Weather-Wind & 40 & {} \\
        \quad Animals-Birds & 32 & {} \\
        \quad Weather-Rain & 23 & {} \\
        \quad Plants-Vegetation & 15 & {} \\
        \quad \textit{Other Natural Subcategories} & 66 & {} \\
        
        \addlinespace
        
        \textbf{Music \& Instruments} & \textbf{94} & \textbf{19.3} \\
        \quad Musical Compositions & 30 & {} \\
        \quad Keyboard-Piano & 17 & {} \\
        \quad String-Guitar & 13 & {} \\
        \quad String-Violin & 12 & {} \\
        \quad \textit{Other Music Subcategories} & 22 & {} \\
        
        \addlinespace
        
        \textbf{Daily Life} & \textbf{50} & \textbf{10.2} \\
        \quad Office-Writing/Typing & 10 & {} \\
        \quad Tools-Power Tools & 8 & {} \\
        \quad \textit{Other Daily Life Subcategories} & 32 & {} \\
        
        \addlinespace
        
        \textbf{Human Voices} & \textbf{37} & \textbf{7.6} \\
        \quad Adults-Conversation & 15 & {} \\
        \quad \textit{Other Human Voices Subcategories} & 22 & {} \\
        
        \addlinespace
        
        % --- For categories with fewer items, listing all might be acceptable ---
        % --- or just list the main category if subcategories are not diverse. ---
        
        \textbf{Transportation} & \textbf{22} & \textbf{4.5} \\
        \quad Vehicles-Cars & 19 & {} \\
        
        \addlinespace
        
        \textbf{Industrial \& Urban} & \textbf{19} & \textbf{3.9} \\
        \quad Fire \& Combustion & 14 & {} \\

        \addlinespace
        
        \textbf{Weapons \& Explosions} & \textbf{12} & \textbf{2.5} \\

        \addlinespace

        \textbf{Special Effects} & \textbf{11} & \textbf{2.3} \\

        \midrule
        
        \textbf{Total Classified} & \textbf{421} & \textbf{84.2} \\
        Unclassified & 79 & 15.8 \\
        \midrule
        \textbf{Grand Total} & \textbf{500} & \textbf{100.0} \\
        
        \bottomrule
    \end{tabularx}
\end{table}
%%%%%%%%%%%%%%%%%%%%%%%%%%%%%%%%%%%%%%%%%%%%%%%%%%%%

%%%%%%%%%%%%%%%%%%%%%%%%%%%%%%%%%%%%%%%%%%%%%%%%%%%%
\begin{figure*}[h!]
    \centering
    \includegraphics[width=1.0\linewidth]{figures/audio_statistics_all.pdf}
    % \vspace{-3mm}
    \caption{Statistical results of set1 and 2.}
    \label{fig:appendix_unibench_statistics_all}
    % \vspace{-5mm}
\end{figure*}
%%%%%%%%%%%%%%%%%%%%%%%%%%%%%%%%%%%%%%%%%%%%%%%%%%%%

%%%%%%%%%%%%%%%%%%%%%%%%%%%%%%%%%%%%%%%%%%%%%%%%%%%%
\begin{figure*}[h!]
    \centering
    \includegraphics[width=1.0\linewidth]{figures/audio_statistics_set1.pdf}
    % \vspace{-3mm}
    \caption{Statistical results of set1.}
    \label{fig:appendix_unibench_statistics_set1}
    % \vspace{-5mm}
\end{figure*}
%%%%%%%%%%%%%%%%%%%%%%%%%%%%%%%%%%%%%%%%%%%%%%%%%%%%

%%%%%%%%%%%%%%%%%%%%%%%%%%%%%%%%%%%%%%%%%%%%%%%%%%%%
\begin{figure*}[h!]
    \centering
    \includegraphics[width=1.0\textwidth]{figures/audio_statistics_set2.pdf}
    % \vspace{-3mm}
    \caption{Statistical results of set2.}
    \label{fig:appendix_unibench_statistics_set2}
\end{figure*}
%%%%%%%%%%%%%%%%%%%%%%%%%%%%%%%%%%%%%%%%%%%%%%%%%%%%



\end{document}